%% kapa: exam=1 round=30 year=2019 period=1 form=A subject=real_estate subject_ko=부동산학원론 seq=8 number=88
%% source: Q-Net 감정평가사 기출문제 (cst00309) artlSeq=5205205 file=extracted/1008/5205205/2019년 제30회 감정평가사 1차시험 문제지/2019년 제30회 감정평가사 1차시험 문제지/2019년도 제30회 감정평가사 1차 1교시 A형.hwp
%% conversion: kordoc 4.24.51 -> tools/gen1.py (auto); review=proofread-sonnet; flags=none
\begin{question}{88}
부동산 수익률에 관한 설명으로 옳지 않은 것을 모두 고른 것은?
\begin{framedtext}
ㄱ. 요구수익률이란 투자자가 투자하기 위한 최대한의 수익률을 말하는 것으로 시간에 대한 비용은 고려하지 않는다. \par
ㄴ. 실현수익률이란 투자가 이루어지고 난 후 현실적으로 달성된 수익률로서 역사적 수익률을 의미한다. \par
ㄷ. 기대수익률이 요구수익률보다 높으면, 대상부동산에 대하여 수요가 증가하여 기대수익률이 상승한다.
\end{framedtext}
\begin{choices}
  \item ㄱ
  \item ㄷ
  \item ㄱ, ㄴ
  \item ㄱ, ㄷ
  \item ㄱ, ㄴ, ㄷ
\end{choices}
\end{question}
\answer{4}
